\section{Introduction}\label{sec:intro}

Software for galactic dynamics has become a standard part of the toolkit
of Galactic astronomy. Codes such as the NEMO Stellar Dynamics Toolbox
\citep{Teuben1995}, \galpy\ \citep{Bovy2015}, \texttt{gala}
\citep{PriceWhelan2017}, and \texttt{AGAMA} \citep{Vasiliev2019}
provide gravitational potentials, orbit integration, action--angle
coordinates, and distribution functions (DFs) in a form that can be
applied directly to observed stars, and they are used both to place
individual objects on their orbits and to build dynamical models that
are fit to data. The reason that such codes are useful is simple: the
motion of a star in a given gravitational potential is a solved
problem, but one that requires careful numerical work, and a
well-tested implementation that is shared saves every user from
repeating that work and from the mistakes made along the way. The
kinematic data from the Gaia mission
\citep{Gaia2016,Gaia2018,Gaia2023}
have made such codes indispensable, because for the first time orbits
can be computed for hundreds of millions of stars.

\citet{Bovy2015} introduced \galpy\ as a \python\ package, with \cc\
implementations of the time-critical parts, for computing orbits,
action--angle coordinates, and DFs in general galactic potentials. It has
since been used in more than a thousand papers (1,056 at the time of
writing, from four in 2012 to about 170 per year since 2024). Most of
these use \galpy\ as a tool: orbits are integrated in a model of the
Milky Way to place objects in the Galaxy and to characterize their
motion. In this way, \galpy\ has been applied to stars in the disk with
measured abundances and ages \citep[\eg,][]{Frankel2020,Xiang2022},
to the substructure of the stellar halo \citep[\eg,][]{Belokurov2020,
Lane2023}, to globular clusters and dwarf galaxies \citep[\eg,][]{Fritz2018,
Pace2022}, to hypervelocity and runaway stars and to the natal kicks of
pulsars and compact objects \citep[\eg,][]{Boubert2018,Marchetti2022,
ODoherty2023}, to open clusters and the star-forming regions near the
Sun \citep[\eg,][]{Zucker2022,Swiggum2024}, to the hosts of exoplanets
and to interstellar objects \citep[\eg,][]{Stefansson2023,Forbes2025},
and to galaxies other than our own \citep[\eg,][]{Tonnesen2019}. A
smaller set of papers, perhaps one in six, uses \galpy\ as a modeling
framework in which a DF or a potential is fit to data: models of tidal
streams and of their perturbations by dark-matter subhalos
\citep[\eg,][]{Bovy2016,Malhan2019,Banik2021}, action-based models of
the disk \citep[\eg,][]{Trick2016,Mackereth2019}, DF and Jeans models
of the halo \citep[\eg,][]{Lane2023}, models of the Galactic bar
\citep[\eg,][]{Hunt2018,Lucey2023}, and measurements of the Milky Way's
mass and potential \citep[\eg,][]{Cautun2020}. It is for this second
kind of use that derivatives matter, because fitting a model to data
requires the likelihood to be evaluated many times over the model's
parameters, and the efficiency of this search is set by whether its
derivatives are available.

\galpy\ has changed much since 2015, but the trends are simple to
state. The number of potentials grew from 22 to 72, with the growth
shifting from specific density profiles to general solvers (basis-function,
multipole, and interpolated potentials) and to wrappers that make any
potential time-dependent, rotating, or tilted, and \galpy\ gained
dissipative forces, such as dynamical friction, and non-inertial
reference frames. The
\orb\ class was rewritten as an arbitrarily-shaped array of orbits that
are integrated in parallel, such that millions of orbits can be handled
at once. DFs for tidal streams were extended to stream gaps
\citep{Sanders2016} and to particle-spray models \citep{Fardal2015,
Chen2025}, and a general framework for spherical DFs appeared. The
action--angle module gained inverse transformations and the orbit module
gained chaos diagnostics, surfaces of section and Lyapunov exponents. In
2026, we implemented the three-dimensional variational equations that
propagate the derivatives of an orbit with respect to its initial
conditions in \cc\ for most of the potential library. Over the same
period the code grew from 36,000 to 93,000 lines of \python\ and from
8,000 to 26,000 lines of \cc, and the test suite grew from 500 to 2,300
tests that cover all of the code.

Over the same decade, numerical computing has been reshaped by
frameworks built for machine learning, chiefly \jax\
\citep{Bradbury2018} and \torch\ \citep{Paszke2019}. These provide
array programming that runs on CPUs and GPUs alike, just-in-time
compilation of whole programs, and automatic differentiation (AD) of
arbitrary code \citep[see][for a review]{Baydin2018}. A dynamics code
written in such a framework gains three things: it runs on a GPU, where
large batches of orbits or stars are processed at once; it can be
compiled into a single program rather than executed one array operation
at a time; and every quantity that it computes---a force, an orbit, an
action, the value of a DF---can be differentiated with respect to every
input, be it an initial condition or a parameter of the potential.
Derivatives are what gradient-based inference requires, whether this is
Hamiltonian Monte Carlo \citep{Duane1987,HoffmanGelman2014}, variational
inference, or the optimization of models with many parameters, such as
neural networks that stand in for a potential or a DF. Codes built on
these frameworks are now in use in cosmology
\citep[\eg,][]{Modi2021,Campagne2023,Li2024} and, over the past two
years, in galactic dynamics. \texttt{galax} \citep{Starkman2026}, a
\jax\ code modeled on \texttt{gala}, provides the halo potentials whose
acceleration directions are compared with the curvature of stellar
streams around other galaxies to constrain the shapes of their halos
\citep{Wu2026b,Euclid2026}, and the orbits used to forecast how well
such curvature tests separate dark-matter models \citep{Starkman2026b}.
\citet{AnauMontel2026} integrate streams perturbed by subhalos with
\texttt{galax} on a GPU and obtain Fisher forecasts by differentiating
the full simulation. Other \jax\ codes fit the progenitor and the host
potential of a mock stream by gradient descent through an $N$-body
simulation \citep{Viterbo2025}, generate the many particle-spray models
needed to infer the distribution of halo flattenings from a population
of streams \citep{Chemaly2026}, and compute the response of streams to
subhalos by integrating linearized equations of motion
\citep{Nibauer2025}. Gradient-based samplers in \jax\ fit a spline
model of the track of GD-1 to measure the accelerations along it
\citep{Nibauer2025b}, and a flexible model of vertical orbits to the
variation of stellar abundances \citep{PriceWhelan2025}. Normalizing
flows for the DF, differentiated automatically, recover the potential
from phase-space snapshots of simulated galaxies
\citep{Green2023,An2021,Buckley2023} and of 5.6 million \textit{Gaia}
stars \citep{Kalda2025}. Differentiating an entire simulation
\citep{Viterbo2025,AnauMontel2026}, or fitting a neural-network
potential to millions of stars \citep{Kalda2025}, is not practical with
codes that only return values.

In this paper, we describe an extension of \galpy\ that adds \jax\ and
\torch\ as \emph{backends} to its existing \numpy\ and \cc\
implementation. The principle that guided this work is that nothing
should change for the existing user: the \numpy\ path produces results
that are byte-for-byte identical to those of the previous version, the
interface is unchanged, and \galpy's existing test suite now runs under
every backend, in eager and in compiled mode, with the same tolerances.
Backends are selected by the type of the arrays that are passed in or
through a context manager; \galpy's code is written against a
backend-agnostic array namespace; and \galpy\ does not compile anything
itself, but is written such that it can be compiled and differentiated
as part of the user's own program. Most of \galpy\ is covered: the
potentials and their forces and derivatives, orbit integration either
inside the backend or through the \cc\ integrators with their variational
equations, the action--angle transformations and their inverses, and the
DFs, including those for tidal streams. All of these are differentiable
with respect to their inputs and parameters, and they run on GPUs. We
list what is not covered where it comes up; the main gaps are the
torus-mapping code and the use of \astropy\ units inside compiled
programs.

Claude, an AI coding agent developed by Anthropic, wrote essentially all
of the code described in this paper, working in the \texttt{Claude Code}
environment under the direction and review of the human author (JB)
between June and October 2026. We discuss how this was done, how we
organized the work, and how we established the correctness of the code
in \sectionname~\ref{sec:process}; a statement on the use of AI, which
also explains the use of ``we'' in a paper with a single human author,
appears at the end of the paper.

The outline of this paper is as follows. \sectionname~\ref{sec:design}
describes the design of the backend framework: how a backend is chosen,
how we made \galpy's code backend-agnostic, the utilities that this
required, and how we made the test suite run under every backend. The
following sections go through \galpy's main modules in turn: potentials
in \sectionname~\ref{sec:potential}, orbit integration and its derivatives
in \sectionname~\ref{sec:orbit}, action--angle coordinates in
\sectionname~\ref{sec:aa}, DFs in \sectionname~\ref{sec:df}, and models
of tidal streams in \sectionname~\ref{sec:streams}. Each of these
sections starts from the general approach and then gives the details that
are specific to the module, with examples of use. \sectionname~\ref{sec:process}
describes the development process. \sectionname~\ref{sec:examples}
presents fuller examples that combine several parts of the code in a
single analysis, and \sectionname~\ref{sec:conclusion} concludes. An
Appendix gives the variational equations behind the orbit derivatives.

Code to reproduce all of the figures in this paper, as well as the
working notes of the development, is available in a \git\ repository
associated with this paper, separate from the \galpy\ repository, at
\\ \centerline{\url{https://github.com/jobovy/galpy-jax}~.}
